\documentclass[a4paper,10pt]{article}

\usepackage[utf8]{inputenc}
\usepackage[T1]{fontenc}
\usepackage[english]{babel}
\usepackage{geometry}
\usepackage{parskip}
\usepackage{hyperref}
\usepackage{titlesec}

\geometry{top=1.0cm, bottom=1.0cm, left=1.0cm, right=1.0cm}
\pagestyle{empty}

\hypersetup{
    pdftitle={Pedro Palacio - Resume},
    pdfauthor={Pedro Palacio},
    colorlinks=true,
    linkcolor=black,
    urlcolor=black,
    citecolor=black,
    bookmarksdepth=1 
}

\setcounter{secnumdepth}{0}

\titleformat{\section}
{\Large\bfseries}
{}
{0em}
{}
[\titlerule\vspace{0.5ex}]

\begin{document}

\begin{center}
    {\LARGE \textbf{Pedro Palacio Bahniuk}} 
    \\ [0.1cm]
    Mandaguari, Paraná, Brazil
    {\textbullet}
    Email: \href{mailto:pedropbhk@gmail.com}{pedropbhk@gmail.com} 
    {\textbullet}
    \href{https://www.linkedin.com/in/pedro-palacio}{linkedin.com/in/pedro-palacio} 
    {\textbullet}
    \href{https://github.com/palaciopedro}{github.com/palaciopedro}
\end{center}

\section{Education}

\subsection*{\textbf{Universidade Estadual de Maringá} \hfill Maringá, Paraná, Brazil}
\textit{B.S. in Software Engineering \hfill Mar 2026 -- Mar 2030}

\section{Professional Summary}

Software Engineering student at Universidade Estadual de Maringá with strong foundations in Python and version control using Git. Experience developing simple applications and responsive web interfaces using HTML and CSS. Currently expanding knowledge in Java and backend development. Seeking an internship opportunity to apply technical skills and gain hands-on industry experience.

\section{Technical Skills}

\begin{itemize}
    \item \textbf{Programming:} Python, Java
    \item \textbf{Web Development:} HTML5, CSS3
    \item \textbf{Database:} PostgreSQL
    \item \textbf{Tools:} Git, GitHub
    \item \textbf{Languages:} Portuguese (Native), English (Advanced)
\end{itemize}

\section{Projects}

\textbf{Personal Website - HTML \& CSS} \\
Developed a responsive website using HTML and CSS to showcase professional information and contact links. Implemented clean layout design and mobile responsiveness. Deployed using GitHub Pages.

\textbf{Task Tracker CLI – Python} \\
Developed a command-line task management application in Python supporting CRUD operations (create, read, update, delete) and task status management. Implemented data persistence using a local JSON file and handled CLI arguments without external libraries. Structured the code in a modular way to improve maintainability

\textbf{Connect API – Java, Spring Boot, MongoDB} \\
Development of a REST API focused on interaction between users and publications, allowing user registration, post creation, and comment addition. The application was built with Spring Boot and Spring Data MongoDB, using NoSQL modeling with embedded documents for comments and references between users and posts. It implements advanced textual queries with regex, filters by date range and search in nested fields, in addition to a layered architecture, DTOs for optimizing data traffic, and CRUD endpoints in JSON, with project management via Maven.

\textbf{Spreadsheet Automation + API Integration - Python, Pandas \& OpenPyXL} \\
Development of a Python application to automate the process of updating Excel spreadsheets, integrating with an API that delivers updated data from the Brazilian Championship Table, eliminating repetitive manual tasks and reducing the risk of errors in data manipulation. The project processes data through a main service and directly updates an existing tab, preserving the original file structure without recreating it. It includes file existence validation, modular code organization, and executable generation for distribution without the need for an installed Python environment.

\textbf{Sales Management API – Java, Spring Boot, Spring Data JPA} \\
Development of a REST API in Java using Spring Boot for managing users, orders, and products. The project uses Spring Data JPA and Hibernate for data persistence, including entity relationship modeling and composite keys for order items. It provides CRUD endpoints using JSON, uses H2 database for testing, supports PostgreSQL, and follows a layered architecture (Controller, Service, Repository) with Maven for dependency management and build.

\end{document}